\parttitle[f-kinetics]{Enzyme Kinetics Equations}{Michaelis-Menten and every extension, with the graphs and the worked numbers.}

\section{Rate, rate constants, order of reaction, and the integrated rate equations biology uses}

Kinetics is the study of how fast a reaction goes, and it answers a different question
from thermodynamics. Thermodynamics tells you whether \ce{C6H12O6 + 6O2 -> 6CO2 + 6H2O}
can happen ($\Delta G^{\circ\prime} = -2870\ \mathrm{kJ\,mol^{-1}}$, so overwhelmingly yes);
kinetics tells you that a bowl of glucose sitting on a bench at $25\degc$ will take
longer than the age of the universe to do it, because the activation barrier is roughly
$\Delta G^{\ddagger} \approx 100\text{--}130\ \mathrm{kJ\,mol^{-1}}$. Every enzyme in
this volume is a device for lowering that barrier without touching $\Delta G$ at all.

\begin{defbox}[title={DEFINITIONS: the vocabulary of rate}]
\nr{f08-kinetics-rate-def}{Reaction rate}
\tm{Rate} (or velocity, $v$) is the change in concentration of a species per unit time,
signed so that it is positive: $v=-d[\mathrm{S}]/dt=+d[\mathrm{P}]/dt$ for
\ce{S -> P}. Standard units in enzymology are
$\mathrm{mol\,L^{-1}\,s^{-1}}$, but assays are usually reported as
$\mu\mathrm{mol\,min^{-1}}$ of product in a stated volume.
\nr{f08-kinetics-rate-law}{Rate law}
The \tm{rate law} is the experimentally determined equation linking $v$ to
concentrations, e.g. $v=k[\mathrm{A}][\mathrm{B}]$. It cannot be read off the balanced
equation; it must be measured.
\nr{f08-kinetics-rate-constant}{Rate constant}
The \tm{rate constant} $k$ is the proportionality factor in the rate law. It depends on
temperature, pH, ionic strength and solvent, but not on concentration.
\nr{f08-kinetics-order-def}{Order of reaction}
The \tm{order} with respect to a species is the exponent on that species in the rate
law. The \tm{overall order} is the sum of the exponents.
\nr{f08-kinetics-molecularity}{Molecularity}
\tm{Molecularity} is the number of molecules colliding in one elementary step; it is a
mechanistic integer. Order is empirical and may be zero, fractional or negative.
\nr{f08-kinetics-initial-rate}{Initial rate}
The \tm{initial rate} $v_0$ is the rate at $t=0$, measured before more than about
$5$ to $10\%$ of substrate is consumed and before product accumulates.
\end{defbox}

The units of $k$ are fixed by the overall order, and this is the single fastest way to
catch an algebra error on an exam. Since $v$ always has units of
$\mathrm{M\,s^{-1}}$, a rate law of overall order $n$ forces $k$ to carry
$\mathrm{M^{1-n}\,s^{-1}}$.

\begin{keybox}[title={MEMORIZE: orders, integrated forms, half-lives and units of $k$}]
\nr{f08-kinetics-integrated-forms}{Integrated rate equations}
For \ce{A -> P} with $[\mathrm{A}]_0$ the starting concentration:
\[
\text{zero order: } v=k,\qquad [\mathrm{A}]_t=[\mathrm{A}]_0-kt,\qquad
t_{1/2}=\frac{[\mathrm{A}]_0}{2k}
\]
\[
\text{first order: } v=k[\mathrm{A}],\qquad \ln[\mathrm{A}]_t=\ln[\mathrm{A}]_0-kt,
\qquad t_{1/2}=\frac{\ln 2}{k}=\frac{0.693}{k}
\]
\[
\text{second order: } v=k[\mathrm{A}]^2,\qquad
\frac{1}{[\mathrm{A}]_t}=\frac{1}{[\mathrm{A}]_0}+kt,\qquad
t_{1/2}=\frac{1}{k[\mathrm{A}]_0}
\]
\nr{f08-kinetics-units-of-k}{Units of the rate constant by order}
Units of $k$: zero order $\mathrm{M\,s^{-1}}$; first order $\mathrm{s^{-1}}$; second
order $\mathrm{M^{-1}\,s^{-1}}$.
\nr{f08-kinetics-halflife-independence}{Concentration independence of first-order half-life}
Only the first-order half-life is independent of $[\mathrm{A}]_0$. That is why
radioactive decay and single-molecule unfolding have fixed half-lives and why
zero-order elimination does not.
\end{keybox}

\begin{tabular}{@{}lllp{3.4cm}@{}}
\toprule
\textbf{Order} & \textbf{Linear plot} & \textbf{Units of $k$} & \textbf{Biological example}\\
\midrule
0 & $[\mathrm{A}]$ vs $t$ & $\mathrm{M\,s^{-1}}$ & Ethanol clearance; any enzyme at $[\mathrm{S}]\gg K_m$\\
1 & $\ln[\mathrm{A}]$ vs $t$ & $\mathrm{s^{-1}}$ & Isotope decay; most drug clearance; $\mathrm{ES\rightarrow E+P}$\\
2 & $1/[\mathrm{A}]$ vs $t$ & $\mathrm{M^{-1}\,s^{-1}}$ & Dimerisation; $\mathrm{E}+\mathrm{S}$ association step\\
Mixed & none & mixed & Michaelis-Menten overall, order $1\rightarrow 0$ in $[\mathrm{S}]$\\
\bottomrule
\end{tabular}

\begin{thmbox}[title={PRINCIPLE: enzymes change the rate constant, never the equilibrium}]
\nr{f08-kinetics-catalyst-principle}{Catalyst principle}
A catalyst lowers $\Delta G^{\ddagger}$ for the forward and reverse directions by the
same amount, so it multiplies $k_f$ and $k_r$ by the same factor and leaves
$K_{eq}=k_f/k_r$ untouched. An enzyme therefore accelerates the approach to
equilibrium and shifts nothing.
\nr{f08-kinetics-eyring}{Eyring equation and rate enhancement}
The transition-state relation
$k=\dfrac{k_BT}{h}e^{-\Delta G^{\ddagger}/RT}$ gives the useful rule of thumb that at
$25\degc$ each $5.7\ \mathrm{kJ\,mol^{-1}}$ removed from the barrier multiplies the
rate by $10$. Enzymes achieve rate enhancements of $10^{6}$ to $10^{17}$, which
corresponds to removing $34$ to $97\ \mathrm{kJ\,mol^{-1}}$. Orotidine
$5'$-monophosphate decarboxylase sits at the top of that range,
$\approx 10^{17}$-fold, with an uncatalysed half-life of $78$ million years.
\end{thmbox}

\begin{warnbox}[title={TRAP: order read off the stoichiometric equation}]
\nr{f08-kinetics-trap-order-stoich}{Order is not stoichiometry}
The wrong belief: for \ce{2A -> P} the rate law must be $v=k[\mathrm{A}]^2$. The
correction: order is empirical. The decomposition of \ce{H2O2} catalysed by iodide is
first order in \ce{H2O2} and first order in \ce{I-} despite the balanced equation
\ce{2H2O2 -> 2H2O + O2}. Only for a single elementary step does molecularity equal
order. In enzymology the overall reaction \ce{S -> P} is first order at low
$[\mathrm{S}]$ and zero order at high $[\mathrm{S}]$ with no change of stoichiometry
whatsoever.
\end{warnbox}

\begin{calcbox}[title={WORKED: first-order decay of a labelled substrate}]
\nr{f08-kinetics-calc-first-order}{Worked first-order decay}
An esterase hydrolyses a substrate present at $[\mathrm{S}]_0=2.0\ \mu\mathrm{M}$,
far below $K_m=250\ \mu\mathrm{M}$, so the loss is first order with observed constant
$k_{obs}=0.0231\ \mathrm{s^{-1}}$.

\ttitle{Half-life.}
\[
t_{1/2}=\frac{0.693}{0.0231\ \mathrm{s^{-1}}}=30.0\ \mathrm{s}
\]

\ttitle{Time to fall to $10\%$ of the start.}
\[
t=\frac{\ln([\mathrm{S}]_0/[\mathrm{S}]_t)}{k_{obs}}
=\frac{\ln 10}{0.0231}=\frac{2.303}{0.0231}=99.7\ \mathrm{s}
\]

\ttitle{Concentration remaining at $t=45\ \mathrm{s}$.}
\[
[\mathrm{S}]_{45}=2.0\,e^{-0.0231\times 45}=2.0\,e^{-1.040}=2.0\times 0.3535
=0.71\ \mu\mathrm{M}
\]

\ttitle{Check.} $99.7\ \mathrm{s}$ is $3.32$ half-lives, and
$2^{-3.32}=0.100$, so $10\%$ remains. $45\ \mathrm{s}$ is $1.5$ half-lives and
$2^{-1.5}=0.354$, giving $0.71\ \mu\mathrm{M}$. Both routes agree. \tm{$\surd$}
\end{calcbox}

\begin{obsbox}[title={WORTH NOTICING: the isotope half-lives an exam expects}]
\nr{f08-kinetics-isotope-halflives}{Half-lives of tracer isotopes}
Every tracer used in classical biology obeys the first-order law above, so
$k=0.693/t_{1/2}$ converts a half-life into a decay constant directly.
$^{32}\mathrm{P}$: $14.3\ \mathrm{d}$. $^{35}\mathrm{S}$: $87.4\ \mathrm{d}$.
$^{125}\mathrm{I}$: $59.4\ \mathrm{d}$. $^{3}\mathrm{H}$: $12.3\ \mathrm{yr}$.
$^{14}\mathrm{C}$: $5730\ \mathrm{yr}$, so
$k=0.693/5730=1.21\times10^{-4}\ \mathrm{yr^{-1}}$. That last number is the one
radiocarbon dating questions want.
\end{obsbox}

\begin{tipbox}[title={TECHNIQUE: deciding the order from a data table in thirty seconds}]
\nr{f08-kinetics-technique-order-id}{Identifying order from data}
Trigger: a table of $[\mathrm{A}]$ against $t$, or of $v$ against $[\mathrm{A}]$.
\begin{itemize}
\item Equal absolute drops in $[\mathrm{A}]$ per equal time interval: zero order.
\item Equal fractional drops (always halving in the same time): first order.
\item $[\mathrm{A}]$ falling more and more slowly than exponentially, with
$1/[\mathrm{A}]$ rising linearly: second order.
\item Doubling $[\mathrm{A}]$ doubles $v$: first order in A. Doubling $[\mathrm{A}]$
quadruples $v$: second order. Doubling $[\mathrm{A}]$ does nothing: zero order,
which for an enzyme means you are already saturated.
\end{itemize}
\end{tipbox}

Zero-order kinetics is the signature of a saturated catalyst, and the clinical example
is worth memorising because the numbers are asked for. Ethanol is cleared by alcohol
dehydrogenase with $K_m\approx 1\ \mathrm{mM}$, while a legally intoxicating blood level
is $17$ to $22\ \mathrm{mM}$, twenty times $K_m$. The enzyme is therefore saturated and
elimination is zero order at roughly $7\ \mathrm{g\,h^{-1}}$ for a $70\ \mathrm{kg}$
adult, quoted in the range $6$ to $10\ \mathrm{g\,h^{-1}}$, equivalent to a fall in
blood alcohol concentration of $0.015$ to $0.018\ \mathrm{g\,per\,100\ mL}$ per hour.
No matter how drunk the person is, the removal rate per hour is the same, which is
exactly what a zero-order law predicts and what a first-order law would not.

\begin{obsbox}[title={WORTH NOTICING: pseudo-first-order is how bimolecular steps are measured}]
\nr{f08-kinetics-pseudo-first-order}{Pseudo-first-order conditions}
A genuinely second-order step $v=k[\mathrm{A}][\mathrm{B}]$ becomes first order in A
if B is held in vast excess, because $[\mathrm{B}]$ then barely changes:
$v=k_{obs}[\mathrm{A}]$ with $k_{obs}=k[\mathrm{B}]$. Plotting $k_{obs}$ against
$[\mathrm{B}]$ gives a straight line through the origin whose slope is the true
second-order constant in $\mathrm{M^{-1}\,s^{-1}}$. Almost every reported
association rate constant in biochemistry, including the $k_1$ of the next section,
was obtained this way.
\end{obsbox}

\section{The Michaelis-Menten equation derived from the steady state assumption, with every symbol defined}

The central observation of enzymology is that $v_0$ against $[\mathrm{S}]$ is not a
straight line but a rectangular hyperbola: it rises steeply at low $[\mathrm{S}]$,
bends over, and approaches a ceiling. Any model that explains this must contain a
saturable binding site, and the minimal such model is one substrate binding reversibly
to one site followed by an irreversible chemical step.

\begin{keybox}[title={MEMORIZE: the minimal mechanism and its four rate constants}]
\nr{f08-kinetics-minimal-mechanism}{Minimal enzyme mechanism}
\[
\ce{E + S <=>[$k_1$][$k_{-1}$] ES ->[$k_2$] E + P}
\]
$k_1$ is the association constant, a second-order constant in
$\mathrm{M^{-1}\,s^{-1}}$, typically $10^{6}$ to $10^{9}$. $k_{-1}$ is the
dissociation constant of the complex, first order, $\mathrm{s^{-1}}$, typically
$1$ to $10^{5}$. $k_2$ is the catalytic step, also written $k_{cat}$, first order,
$\mathrm{s^{-1}}$. The step \ce{E + P -> ES} is neglected because $v_0$ is measured
when $[\mathrm{P}]\approx 0$.
\end{keybox}

\begin{defbox}[title={DEFINITIONS: every symbol in the derivation}]
\nr{f08-kinetics-symbols}{Symbols of the Michaelis-Menten derivation}
\begin{itemize}
\item $[\mathrm{E}]_T$: total enzyme concentration, free plus bound. Constant, because
enzyme is a catalyst. Typical assay value $10^{-11}$ to $10^{-8}\ \mathrm{M}$.
\item $[\mathrm{E}]$: free enzyme concentration at any instant.
\item $[\mathrm{ES}]$: concentration of the enzyme-substrate complex.
\item $[\mathrm{S}]$: free substrate concentration, taken as equal to the amount added
because $[\mathrm{S}]\gg[\mathrm{E}]_T$.
\item $v_0$: initial velocity, $\mathrm{M\,s^{-1}}$.
\item $V_{\max}$: the limiting velocity as $[\mathrm{S}]\rightarrow\infty$,
$\mathrm{M\,s^{-1}}$.
\item $K_m$: the Michaelis constant, units of concentration, $\mathrm{M}$.
\item $f_{ES}=[\mathrm{ES}]/[\mathrm{E}]_T$: fraction of enzyme in complex, which
equals $v_0/V_{\max}$.
\end{itemize}
\end{defbox}

\begin{thmbox}[title={PRINCIPLE: the steady state assumption}]
\nr{f08-kinetics-steady-state}{Steady state assumption}
After a pre-steady-state burst lasting microseconds to milliseconds, $[\mathrm{ES}]$
becomes constant because it is formed as fast as it is destroyed:
\[
\frac{d[\mathrm{ES}]}{dt}=k_1[\mathrm{E}][\mathrm{S}]-k_{-1}[\mathrm{ES}]-k_2[\mathrm{ES}]=0
\]
This is valid whenever $[\mathrm{S}]\gg[\mathrm{E}]_T$, typically by a factor of
$10^{3}$ or more, so that the substrate reservoir is not measurably drained during the
approach to steady state. It does not require $k_2\ll k_{-1}$, which is what
distinguishes it from the earlier equilibrium treatment.
\end{thmbox}

The derivation is four lines and is worth being able to reproduce from memory, because
the Semifinal Exam has asked for it and because every inhibition equation later in this
volume is the same algebra with one extra term.

\begin{thmbox}[title={PRINCIPLE: derivation of the Michaelis-Menten equation}]
\nr{f08-kinetics-mm-derivation}{Michaelis-Menten derivation}
Step 1, steady state: $k_1[\mathrm{E}][\mathrm{S}]=(k_{-1}+k_2)[\mathrm{ES}]$.
Step 2, define $K_m$ as the ratio of the constants that destroy the complex to the one
that makes it:
\[
K_m=\frac{k_{-1}+k_2}{k_1}\quad\Longrightarrow\quad
[\mathrm{E}][\mathrm{S}]=K_m[\mathrm{ES}]
\]
Step 3, conservation of enzyme: $[\mathrm{E}]=[\mathrm{E}]_T-[\mathrm{ES}]$, so
$([\mathrm{E}]_T-[\mathrm{ES}])[\mathrm{S}]=K_m[\mathrm{ES}]$, and solving,
\[
[\mathrm{ES}]=\frac{[\mathrm{E}]_T[\mathrm{S}]}{K_m+[\mathrm{S}]}
\]
Step 4, the observed rate is the rate of the catalytic step, $v_0=k_2[\mathrm{ES}]$,
and the maximum possible rate is reached when all enzyme is complexed,
$V_{\max}=k_2[\mathrm{E}]_T$. Substituting:
\[
v_0=\frac{V_{\max}[\mathrm{S}]}{K_m+[\mathrm{S}]}
\]
\end{thmbox}

\begin{keybox}[title={MEMORIZE: the Michaelis-Menten equation and its three limits}]
\nr{f08-kinetics-mm-equation}{Michaelis-Menten equation}
\[
v_0=\frac{V_{\max}[\mathrm{S}]}{K_m+[\mathrm{S}]}
\]
\nr{f08-kinetics-mm-limits}{Limiting behaviour of the Michaelis-Menten equation}
\begin{itemize}
\item $[\mathrm{S}]\ll K_m$: the $K_m$ term dominates the denominator and
$v_0\approx(V_{\max}/K_m)[\mathrm{S}]$, first order in substrate. The slope
$V_{\max}/K_m$ equals $(k_{cat}/K_m)[\mathrm{E}]_T$.
\item $[\mathrm{S}]=K_m$: $v_0=V_{\max}/2$ exactly. This is the operational definition
of $K_m$.
\item $[\mathrm{S}]\gg K_m$: $v_0\approx V_{\max}$, zero order in substrate. The enzyme
is saturated and adding more substrate does nothing.
\end{itemize}
\nr{f08-kinetics-mm-conditions}{Conditions for the Michaelis-Menten equation to hold}
Conditions: single substrate, one binding site per active site, no cooperativity,
initial-rate measurement with $[\mathrm{P}]\approx 0$, no product inhibition,
$[\mathrm{S}]\gg[\mathrm{E}]_T$, constant pH and temperature, and steady state
established.
\end{keybox}

\begin{expbox}[title={EXPERIMENT: Michaelis and Menten on invertase, 1913}]
\nr{f08-kinetics-exp-michaelis-menten}{Michaelis and Menten 1913}
\ttitle{Design.} Leonor Michaelis and Maud Menten followed the hydrolysis of sucrose to
glucose plus fructose by yeast invertase, using the change in optical rotation. The
methodological advances were what mattered: they buffered the reaction, which fixed the
pH that Victor Henri in $1903$ had not controlled; they measured only the initial rate,
which removed product inhibition by fructose; and they varied $[\mathrm{S}]$ over
about a $100$-fold range.

\ttitle{Result.} $v_0$ against $[\mathrm{S}]$ was a rectangular hyperbola. Fitting gave
a half-saturation constant of about $0.017\ \mathrm{M}$ sucrose.

\ttitle{What it established.} That the enzyme forms a discrete, saturable complex with
its substrate, and that one constant with units of concentration plus one limiting
velocity describe the whole curve. Henri had written the equation first, in $1903$, but
without buffering or initial-rate measurement he could not test it; the equation is
occasionally called Henri-Michaelis-Menten for that reason.
\end{expbox}

\begin{expbox}[title={EXPERIMENT: Briggs and Haldane, 1925, and why the assumption changed}]
\nr{f08-kinetics-exp-briggs-haldane}{Briggs-Haldane steady state, 1925}
\ttitle{Design.} George Briggs and J.B.S. Haldane re-derived the same hyperbola without
assuming that \ce{ES} is in equilibrium with \ce{E + S}. Michaelis and Menten had
required $k_2\ll k_{-1}$, giving $K_m=k_{-1}/k_1=K_S$, a true dissociation constant.
Briggs and Haldane assumed only $d[\mathrm{ES}]/dt=0$.

\ttitle{Result.} The identical equation appears, but with
$K_m=(k_{-1}+k_2)/k_1$.

\ttitle{What it established.} That the hyperbola is compatible with any relative size
of $k_2$ and $k_{-1}$, so fitting a hyperbola tells you nothing about the mechanism
beyond saturation. It also established that $K_m\geq K_S$ always, with equality only
when $k_2\ll k_{-1}$. For acetylcholinesterase, where $k_2$ is very large,
$K_m$ exceeds $K_S$ substantially.
\end{expbox}

\begin{warnbox}[title={TRAP: treating $K_m$ as a dissociation constant}]
\nr{f08-kinetics-trap-km-not-kd}{$K_m$ is not always $K_d$}
The wrong belief: $K_m$ measures how tightly the enzyme binds its substrate, so a low
$K_m$ always means tight binding. The correction: $K_m=(k_{-1}+k_2)/k_1$ contains the
catalytic constant. Only when $k_2\ll k_{-1}$ does $K_m$ collapse to
$K_S=k_{-1}/k_1=K_d$. When $k_2\gg k_{-1}$ the enzyme converts substrate faster than
it releases it, $K_m\approx k_2/k_1$, and $K_m$ is then a measure of catalytic
throughput, not affinity. Safe phrasing on an exam: $K_m$ is the substrate
concentration giving half-maximal velocity, and it is an inverse measure of
\emph{apparent} affinity.
\end{warnbox}

\begin{tabular}{@{}lp{3.0cm}p{4.2cm}@{}}
\toprule
\textbf{Treatment} & \textbf{Assumption} & \textbf{Meaning of the constant}\\
\midrule
Henri 1903 & Rapid equilibrium, unbuffered & Hyperbola written but untested\\
Michaelis-Menten 1913 & $k_2\ll k_{-1}$ & $K_S=k_{-1}/k_1$, a true $K_d$\\
Briggs-Haldane 1925 & $d[\mathrm{ES}]/dt=0$ & $K_m=(k_{-1}+k_2)/k_1\geq K_S$\\
\bottomrule
\end{tabular}

\begin{obsbox}[title={WORTH NOTICING: the same hyperbola appears in four other places}]
\nr{f08-kinetics-hyperbola-analogues}{Hyperbolic saturation analogues}
The form $y=y_{\max}x/(K+x)$ is not special to enzymes. Myoglobin oxygen binding uses
$Y=p\mathrm{O_2}/(P_{50}+p\mathrm{O_2})$ with $P_{50}\approx 2.8\ \mathrm{mmHg}$.
Carrier-mediated facilitated diffusion of glucose by GLUT1 uses
$K_m\approx 1.5\ \mathrm{mM}$ with the same algebra. Receptor occupancy uses $K_d$.
Nutrient-limited microbial growth uses the Monod equation
$\mu=\mu_{\max}S/(K_S+S)$, and predator functional response type II uses Holling's
disc equation. Recognising the shared algebra means one derivation covers all five.
\end{obsbox}

\section{$K_m$, $V_{\max}$, $k_{cat}$, catalytic efficiency, and turnover number, with worked calculations and units}

Four numbers describe a Michaelis-Menten enzyme, and confusing them is the most common
source of lost marks on this topic. Two are properties of the enzyme alone
($k_{cat}$ and $k_{cat}/K_m$), one depends on how much enzyme you happened to put in
the tube ($V_{\max}$), and one is a concentration ($K_m$).

\begin{defbox}[title={DEFINITIONS: the four kinetic parameters}]
\nr{f08-kinetics-km-def}{Michaelis constant $K_m$}
\tm{$K_m$} is the substrate concentration at which $v_0=V_{\max}/2$. Units of
concentration, usually $\mathrm{mM}$ or $\mu\mathrm{M}$. Independent of
$[\mathrm{E}]_T$.
\nr{f08-kinetics-vmax-def}{Maximum velocity $V_{\max}$}
\tm{$V_{\max}$} is the velocity at saturating substrate, $V_{\max}=k_{cat}[\mathrm{E}]_T$.
Units of $\mathrm{M\,s^{-1}}$ or $\mu\mathrm{mol\,min^{-1}}$. Directly proportional to
enzyme concentration, so it is not a constant of the enzyme.
\nr{f08-kinetics-kcat-def}{Catalytic constant $k_{cat}$}
\tm{$k_{cat}$}, the \tm{turnover number}, is the number of substrate molecules
converted per active site per second at saturation. Units $\mathrm{s^{-1}}$.
$k_{cat}=V_{\max}/[\mathrm{E}]_T$. For the minimal mechanism $k_{cat}=k_2$; for longer
mechanisms it is the composite constant of the slowest step after binding.
\nr{f08-kinetics-efficiency-def}{Catalytic efficiency $k_{cat}/K_m$}
\tm{Catalytic efficiency} or the \tm{specificity constant} is $k_{cat}/K_m$, units
$\mathrm{M^{-1}\,s^{-1}}$. It is the apparent second-order rate constant for the
reaction of free enzyme with free substrate at low $[\mathrm{S}]$, and it is the
number that decides which of two competing substrates an enzyme prefers.
\end{defbox}

\begin{keybox}[title={MEMORIZE: the conversions between the four parameters}]
\nr{f08-kinetics-parameter-conversions}{Conversions among kinetic parameters}
\[
V_{\max}=k_{cat}[\mathrm{E}]_T,\qquad
k_{cat}=\frac{V_{\max}}{[\mathrm{E}]_T},\qquad
\frac{V_{\max}}{K_m}=\frac{k_{cat}}{K_m}[\mathrm{E}]_T
\]
\nr{f08-kinetics-subunit-correction}{Active-site correction to $k_{cat}$}
If the enzyme has $n$ active sites per molecule, $[\mathrm{E}]_T$ in these equations
means the concentration of \emph{active sites}, so divide $V_{\max}$ by
$n\times[\text{protein}]$. Aspartate transcarbamoylase has six catalytic sites,
haemoglobin four, and glyceraldehyde-3-phosphate dehydrogenase four; forgetting the
factor gives a $k_{cat}$ wrong by exactly that factor.
\end{keybox}

\begin{calcbox}[title={WORKED: $k_{cat}$ and catalytic efficiency from an assay}]
\nr{f08-kinetics-calc-kcat}{Worked $k_{cat}$ and efficiency}
A purified monomeric dehydrogenase is assayed in $1.0\ \mathrm{mL}$ at $37\degc$. The
enzyme is at $5.0\ \mathrm{nM}$. At saturating substrate the assay produces
$15\ \mathrm{nmol}$ of product per minute. $K_m$ was measured as $0.40\ \mathrm{mM}$.

\ttitle{Moles of enzyme in the cuvette.}
\[
n_E=(5.0\times10^{-9}\ \mathrm{mol\,L^{-1}})\times(1.0\times10^{-3}\ \mathrm{L})
=5.0\times10^{-12}\ \mathrm{mol}=5.0\ \mathrm{pmol}
\]

\ttitle{$k_{cat}$.} Convert the rate to pmol per minute:
$15\ \mathrm{nmol\,min^{-1}}=1.5\times10^{4}\ \mathrm{pmol\,min^{-1}}$.
\[
k_{cat}=\frac{1.5\times10^{4}\ \mathrm{pmol\,min^{-1}}}{5.0\ \mathrm{pmol}}
=3.0\times10^{3}\ \mathrm{min^{-1}}
=\frac{3.0\times10^{3}}{60}=50\ \mathrm{s^{-1}}
\]

\ttitle{$V_{\max}$ as a concentration rate.}
$15\ \mathrm{nmol\,min^{-1}}$ in $1.0\ \mathrm{mL}$ is
$15\ \mu\mathrm{mol\,L^{-1}\,min^{-1}}=0.25\ \mu\mathrm{M\,s^{-1}}$.

\ttitle{Catalytic efficiency.}
\[
\frac{k_{cat}}{K_m}=\frac{50\ \mathrm{s^{-1}}}{4.0\times10^{-4}\ \mathrm{M}}
=1.25\times10^{5}\ \mathrm{M^{-1}\,s^{-1}}
\]

\ttitle{Check.} Cross-check $V_{\max}=k_{cat}[\mathrm{E}]_T
=50\ \mathrm{s^{-1}}\times5.0\times10^{-9}\ \mathrm{M}
=2.5\times10^{-7}\ \mathrm{M\,s^{-1}}=0.25\ \mu\mathrm{M\,s^{-1}}$, which matches the
independently computed value. The efficiency is $10^{4}$-fold below the diffusion
limit of $\approx10^{9}\ \mathrm{M^{-1}\,s^{-1}}$, so this is an ordinary, not a
perfected, enzyme. \tm{$\surd$}
\end{calcbox}

\begin{tabular}{@{}lllll@{}}
\toprule
\textbf{Enzyme} & \textbf{Substrate} & \textbf{$K_m$ (mM)} & \textbf{$k_{cat}$ (s$^{-1}$)} & \textbf{$k_{cat}/K_m$}\\
\midrule
Catalase & \ce{H2O2} & $25$ & $4\times10^{7}$ & $1.6\times10^{9}$\\
Carbonic anhydrase II & \ce{CO2} & $12$ & $1\times10^{6}$ & $8.3\times10^{7}$\\
Acetylcholinesterase & ACh & $0.095$ & $1.4\times10^{4}$ & $1.5\times10^{8}$\\
Fumarase & fumarate & $0.005$ & $800$ & $1.6\times10^{8}$\\
Triose-P isomerase & GAP & $0.47$ & $4.3\times10^{3}$ & $9.1\times10^{6}$\\
Superoxide dismutase & \ce{O2^-} & $0.36$ & $1\times10^{6}$ & $2.8\times10^{9}$\\
Chymotrypsin & Ac-Phe-OEt & $15$ & $0.14$ & $9.3$\\
Hexokinase & glucose & $0.05$ & $200$ & $4\times10^{6}$\\
Glucokinase & glucose & $10$ & $60$ & $6\times10^{3}$\\
\bottomrule
\end{tabular}

Published values for these enzymes vary by up to a factor of five between textbooks,
because they depend on pH, temperature, ionic strength and which isoform was used.
Catalase $K_m$ for \ce{H2O2} is quoted anywhere from $25\ \mathrm{mM}$ to over
$1\ \mathrm{M}$ because the enzyme is inactivated by its own substrate at high
concentration; $k_{cat}$ is quoted from $1\times10^{7}$ to $4\times10^{7}\
\mathrm{s^{-1}}$. Quote a value with the range and you cannot be marked wrong.

\begin{thmbox}[title={PRINCIPLE: the diffusion limit and catalytic perfection}]
\nr{f08-kinetics-diffusion-limit}{Diffusion limit on $k_{cat}/K_m$}
A bimolecular encounter in water cannot happen faster than diffusion allows, which puts
a ceiling of $10^{8}$ to $10^{9}\ \mathrm{M^{-1}\,s^{-1}}$ on $k_{cat}/K_m$ for a
normal-sized substrate.
\nr{f08-kinetics-catalytic-perfection}{Catalytic perfection}
An enzyme whose $k_{cat}/K_m$ reaches that ceiling is \tm{catalytically perfect}: every
encounter with substrate yields product, so no further improvement in the chemistry can
increase the rate. Triose phosphate isomerase, fumarase, catalase, carbonic anhydrase
and superoxide dismutase are the standard examples. Superoxide dismutase is reported
at $2\times10^{9}$ to $7\times10^{9}\ \mathrm{M^{-1}\,s^{-1}}$, apparently above the
limit, because electrostatic steering funnels the anionic substrate into the active
site.
\end{thmbox}

\begin{warnbox}[title={TRAP: $V_{\max}$ quoted as a property of the enzyme}]
\nr{f08-kinetics-trap-vmax-property}{$V_{\max}$ depends on enzyme amount}
The wrong belief: enzyme A has a higher $V_{\max}$ than enzyme B, therefore A is the
better catalyst. The correction: $V_{\max}=k_{cat}[\mathrm{E}]_T$, so doubling the
enzyme doubles $V_{\max}$ while $k_{cat}$, $K_m$ and $k_{cat}/K_m$ are all unchanged.
Comparisons between enzymes must use $k_{cat}$ or $k_{cat}/K_m$. A question that says
"the amount of enzyme was doubled" expects you to say $V_{\max}$ doubles and $K_m$ is
unchanged; the saturation curve rises but the half-saturation point stays put.
\end{warnbox}

\begin{warnbox}[title={TRAP: low $K_m$ equals better enzyme}]
\nr{f08-kinetics-trap-low-km-better}{Low $K_m$ is not the same as better}
The wrong belief: the enzyme with the lower $K_m$ is physiologically superior. The
correction: a low $K_m$ means the enzyme is already near-saturated at physiological
substrate concentration, so it runs flat out and cannot respond to a rise in
substrate. Hexokinase, $K_m\approx 0.05\ \mathrm{mM}$, is saturated at the
$4$ to $5\ \mathrm{mM}$ glucose of blood and keeps muscle supplied at all times.
Glucokinase, $K_m\approx 10\ \mathrm{mM}$ with a sigmoidal response, is deliberately
unsaturated so that hepatic glucose uptake rises steeply when portal glucose climbs
from $5$ to $15\ \mathrm{mM}$ after a meal. Neither enzyme is better; they are tuned
to different jobs. The same logic distinguishes cytosolic
aldehyde dehydrogenase ($K_m\approx 30\ \mu\mathrm{M}$) from the mitochondrial ALDH2
($K_m\approx 0.2\ \mu\mathrm{M}$).
\end{warnbox}

\begin{obsbox}[title={WORTH NOTICING: $K_m$ usually sits near physiological $[\mathrm{S}]$}]
\nr{f08-kinetics-km-near-physiological}{$K_m$ tracks physiological substrate}
For most enzymes $K_m$ lies within about an order of magnitude of the cellular
concentration of the substrate. This is not a coincidence: at $[\mathrm{S}]\approx K_m$
the enzyme runs at half speed, which is the point of maximum sensitivity, because
$dv/d[\mathrm{S}]$ relative to $v$ is largest in the rising part of the hyperbola.
An enzyme operating at $[\mathrm{S}]=10K_m$ would be a fixed-rate pump; one at
$[\mathrm{S}]=0.01K_m$ would waste most of its capacity.
\end{obsbox}

\section{Reading kinetic parameters from a saturation curve, and from initial rate data}

The direct hyperbolic plot is how kinetics is actually done today, because non-linear
regression on the untransformed data is statistically correct. An exam still expects you
to read parameters off a curve and off a table by hand, and the errors that this
introduces are themselves examinable.

\begin{tipbox}[title={TECHNIQUE: extracting $V_{\max}$ and $K_m$ from a saturation curve}]
\nr{f08-kinetics-technique-read-curve}{Reading a saturation curve}
Trigger: a graph of $v_0$ on the vertical axis against $[\mathrm{S}]$ on the horizontal,
curving to a plateau.
\begin{enumerate}
\item Identify the horizontal asymptote. That value is $V_{\max}$. Check that the last
two or three points differ by less than $5\%$; if the curve is still climbing you are
reading a lower bound, not $V_{\max}$.
\item Halve it. Draw a horizontal line at $V_{\max}/2$.
\item Drop a vertical from where that line meets the curve. The $[\mathrm{S}]$ value is
$K_m$, in whatever concentration units the axis carries.
\item Confirm with the initial slope: the tangent at the origin has slope
$V_{\max}/K_m$, so $K_m$ predicted from the slope should match.
\end{enumerate}
\end{tipbox}

\begin{keybox}[title={MEMORIZE: the $81$-fold rule and the saturation table}]
\nr{f08-kinetics-81-fold-rule}{Eighty-one-fold rule}
Going from $10\%$ to $90\%$ of $V_{\max}$ requires an $81$-fold increase in
$[\mathrm{S}]$, from $0.111K_m$ to $9K_m$. This is a property of every rectangular
hyperbola and is the reason a saturation curve is so hard to saturate.
\nr{f08-kinetics-fraction-table}{Fractional velocity at multiples of $K_m$}
\[
\frac{v_0}{V_{\max}}=\frac{[\mathrm{S}]/K_m}{1+[\mathrm{S}]/K_m}
\]
So $[\mathrm{S}]=K_m$ gives exactly $0.500$; $[\mathrm{S}]=10K_m$ gives $0.909$;
$[\mathrm{S}]=100K_m$ gives $0.990$. To be within $1\%$ of $V_{\max}$ you need
$[\mathrm{S}]=99K_m$, which is often experimentally impossible because of solubility.
\end{keybox}

\begin{tabular}{@{}llll@{}}
\toprule
$[\mathrm{S}]/K_m$ & $v_0/V_{\max}$ & $[\mathrm{S}]/K_m$ & $v_0/V_{\max}$\\
\midrule
$0.1$ & $0.091$ & $3$ & $0.750$\\
$0.111$ & $0.100$ & $5$ & $0.833$\\
$0.25$ & $0.200$ & $9$ & $0.900$\\
$0.5$ & $0.333$ & $10$ & $0.909$\\
$1$ & $0.500$ & $19$ & $0.950$\\
$2$ & $0.667$ & $99$ & $0.990$\\
\bottomrule
\end{tabular}

\begin{warnbox}[title={TRAP: eyeballing $V_{\max}$ off an unsaturated curve}]
\nr{f08-kinetics-trap-vmax-eyeball}{Underestimating $V_{\max}$ by eye}
The wrong belief: the highest measured velocity is close enough to $V_{\max}$. The
correction: if the largest $[\mathrm{S}]$ used was only $3K_m$, the highest velocity is
$0.75V_{\max}$, so $V_{\max}$ is underestimated by $25\%$ and $K_m$, read as the
substrate giving half of that wrong ceiling, is underestimated too. Design assays to
span $0.2K_m$ to $5K_m$ at minimum, with at least $6$ to $8$ substrate concentrations
spaced geometrically, not arithmetically. If a graph in a question stops climbing
suspiciously early, suspect substrate inhibition rather than saturation.
\end{warnbox}

\begin{calcbox}[title={WORKED: parameters read straight from an initial-rate table}]
\nr{f08-kinetics-calc-read-table}{Worked reading of an initial-rate table}
An enzyme assay at $30\degc$, pH $7.4$, gives these initial rates.

\ttitle{Data.}
$[\mathrm{S}]$ in $\mathrm{mM}$: $0.10$, $0.25$, $0.50$, $1.00$, $2.00$, $5.00$,
$10.00$, $20.00$.
$v_0$ in $\mu\mathrm{mol\,min^{-1}}$: $6.9$, $15.0$, $25.0$, $37.5$, $50.0$, $62.5$,
$68.2$, $72.7$.

\ttitle{Step 1, bound $V_{\max}$.} The last two rates differ by $6.6\%$, so the curve
is still rising and $72.7$ is a lower bound. The pattern of the numbers suggests a
ceiling near $75$.

\ttitle{Step 2, test the guess.} If $V_{\max}=75.0$ and $K_m=1.00\ \mathrm{mM}$, then
at $[\mathrm{S}]=1.00$, $v_0=75.0\times1.00/(1.00+1.00)=37.5$. Matches exactly.

\ttitle{Step 3, verify at three more points.}
\[
[\mathrm{S}]=0.25:\ \frac{75.0\times0.25}{1.25}=15.0\quad\text{matches}
\]
\[
[\mathrm{S}]=5.00:\ \frac{75.0\times5.00}{6.00}=62.5\quad\text{matches}
\]
\[
[\mathrm{S}]=20.0:\ \frac{75.0\times20.0}{21.0}=71.4\quad\text{vs }72.7\text{ measured}
\]

\ttitle{Step 4, $K_m$ by the half-velocity route.} $V_{\max}/2=37.5$, and the table
shows $v_0=37.5$ at $[\mathrm{S}]=1.00\ \mathrm{mM}$, so
$K_m=1.00\ \mathrm{mM}$ directly.

\ttitle{Check.} The initial slope should be $V_{\max}/K_m=75.0\
\mu\mathrm{mol\,min^{-1}\,mM^{-1}}$. From the lowest point,
$6.9/0.10=69$, which is $92\%$ of $75$; the shortfall is expected because
$[\mathrm{S}]=0.1K_m$ is not quite in the linear region, where the predicted fraction
is $0.091$ rather than $0.100$. Indeed $75.0\times0.091=6.8$, matching the measured
$6.9$. \tm{$\surd$} The residual at $20\ \mathrm{mM}$ hints at a second, weaker
activating site or simply $1.8\%$ measurement scatter.
\end{calcbox}

\begin{tipbox}[title={TECHNIQUE: making sure a rate is an initial rate}]
\nr{f08-kinetics-technique-initial-rate}{Securing a true initial rate}
Trigger: any progress curve, or any question that gives product formed after a fixed
incubation.
\begin{itemize}
\item Take the tangent at $t=0$, not the chord from $0$ to the last point. The chord
always underestimates.
\item Keep total substrate depletion under $10\%$, and preferably under $5\%$. At
$10\%$ depletion an enzyme running at $[\mathrm{S}]=K_m$ has already slowed by about
$5\%$.
\item Keep $[\mathrm{E}]_T$ at least $10^{3}$-fold below $[\mathrm{S}]$ so the steady
state holds and free substrate equals added substrate.
\item Watch for a lag from enzyme activation or from a coupled indicator reaction, and
for a burst from a rapid first turnover.
\item Confirm linearity by showing that $v_0$ is proportional to $[\mathrm{E}]_T$ over
at least a fourfold range. If it is not, something other than the enzyme is limiting,
such as oxygen supply or detector saturation.
\end{itemize}
\end{tipbox}

\begin{tabular}{@{}p{3.1cm}p{3.6cm}p{3.6cm}@{}}
\toprule
\textbf{Artefact in a progress curve} & \textbf{Appearance} & \textbf{Cause and fix}\\
\midrule
Chord instead of tangent & $v_0$ too low, more so at low $[\mathrm{S}]$ & Sample earlier; use the first $5\%$ of turnover\\
Lag phase & Curve concave up at the start & Coupled indicator, cofactor binding, or slow mixing; pre-incubate\\
Burst then slowing & Steep first $100\ \mathrm{ms}$ & Rapid acylation before slow deacylation; needs stopped-flow\\
Plateau too early & False low $V_{\max}$ & Substrate inhibition, product inhibition, or enzyme denaturation\\
Non-proportional to $[\mathrm{E}]_T$ & Doubling enzyme gives $<2\times$ rate & Detector saturation or diffusion-limited supply\\
\bottomrule
\end{tabular}

\begin{obsbox}[title={WORTH NOTICING: why non-linear fitting beat the linear plots}]
\nr{f08-kinetics-nonlinear-preferred}{Non-linear regression on the hyperbola}
Fitting $v_0=V_{\max}[\mathrm{S}]/(K_m+[\mathrm{S}])$ directly by least squares assumes
that the error in $v_0$ is uniform, which is roughly true for a spectrophotometric
assay. Every linear transformation in the next section redistributes that error and
therefore gives a different, usually worse, estimate. The linear plots survive because
they make the \emph{shape} of an inhibition effect visible at a glance, not because
they give better numbers.
\end{obsbox}
